\section*{第 \ref{ch:b3:nuclear-physics} 章——核物理}

\begin{solution}{exo:b3:nuclear-physics:1}
(a) $^{12}$C: $6, 6, 12$; $^{14}$C: $6, 8, 14$; $^{235}$U:
$92, 143, 235$；$^{238}$U：$92, 146, 238$。(b) 碳的两种；
铀的两种。(c) 化学性质取决于由 $Z$ 固定的电子云；
稳定性取决于内部的 $N$--$Z$ 平衡。(d) 同量异位素——
由 $\beta^-$ 衰变相连：$^{14}\text{C} \to {}^{14}\text{N}
+ \text{e}^- + \bar\nu$，即放射性碳钟的“滴答”。
\end{solution}

\begin{solution}{exo:b3:nuclear-physics:2}
(a) $\rho = Am_{\text{u}}/\tfrac43\pi r_0^3A =
3m_{\text{u}}/4\pi r_0^3 \approx \qty{2.3e17}{kg/m^3}$: $A$
约去。(b) $\qty{5}{mL} \to \qty{1.2e12}{kg}$——一茶匙的量就重达十亿
吨。(c) $(R/R_{\text{原子}})^3 \approx
(4.7\times10^{-5})^3 \approx 10^{-13}$：原子基本是空的，
中间嵌着一个很重的小点。(d) 体积 $\propto A$ 表明每个核子
占据固定的空间，且只与紧邻的核子结合：
核力在 \unit{fm} 的作用范围内饱和。
\end{solution}

\begin{solution}{exo:b3:nuclear-physics:3}
(a) $\Delta m = 2(1.007276 + 1.008665) - 4.001506 =
\qty{0.030376}{u}$: $B = \qty{28.3}{MeV}$. (b) $B/A =
\qty{7.07}{MeV}$. (c) $28.3\,\text{MeV}/13.6\,\text{eV} \approx
2\times10^6$：差了六个数量级——这正是核燃料
远超化学燃料约一百万倍。(d) $\alpha$ 粒子是一个
预先组装好、极其廉价的“包裹”：放出它
便以 \qty{28.3}{MeV} 的优惠价输出四个核子，
这是任何单个核子都无法匹敌的。
\end{solution}

\begin{solution}{exo:b3:nuclear-physics:4}
(a) $N = N_{\text{A}}/226 = \qty{2.66e21}{}$. (b) $\lambda =
\ln2/(1600\times\qty{3.156e7}{s}) = \qty{1.37e-11}{s^{-1}}$. (c)
$\mathcal A = \lambda N = \qty{3.7e10}{Bq}$：一居里，
按定义——居里夫妇自己的那一克镭。(d) 三个半衰期后：
$\qty{0.125}{g}$ 的镭（其余已变为氡及其
子体，一路奔向铅）。
\end{solution}

\begin{solution}{exo:b3:nuclear-physics:5}
(a) 体积项 $+882$；表面项 $-261$；库仑项 $-0.711\times676/3.83
= -126$；不对称项 $-23.7\times16/56 = -6.8$（\unit{MeV}）。(b)
考虑配对项 $+12/\sqrt{56} = +1.6$：$B \approx \qty{490}{MeV}$，
$B/A = 8.76$，与实测的 $8.79$ 相比——偏差为千分之三，
由液滴模型得到。(c) 体积项与表面项$+$库仑项相抗衡：它们
增长率的交点正是铁处极大值的成因。
(d) 库仑力是长程的：每个质子都排斥其余所有质子，
$\propto Z^2$，而饱和的强核力却只会
随 $\propto A$ 增长——这便是重核最终的败因。
\end{solution}

\begin{solution}{exo:b3:nuclear-physics:6}
(a) 对库仑项与不对称项作 $\partial/\partial Z$ 运算，
得 $2a_{\text{C}}Z/A^{1/3} = 4a_{\text{A}}(A - 2Z)/A$，解之即可。
(b) $Z^* = 50.5/1.163 = 43.4$：$A = 101$ 时，大自然的选择是
钌，$Z = 44$。(c) 若无库仑项，$Z^* = A/2$：两个费米
阶梯（见\cref{ch:b3:quantum-statistics}），最省的做法是
将两者填到同一高度。(d) 碎片在谷底的位置为 $Z^* \approx 41$：
缺五个质子——五次 $\beta^-$ 衰变的级联
可将其爬回谷底，每次把一个中子变为质子，并放出
一个电子和一个反中微子。
\end{solution}

\begin{solution}{exo:b3:nuclear-physics:7}
(a) $t = T_{1/2}\ln(1/0.3)/\ln2 = 5730\times1.74 \approx
\qty{e4}{}$ 年：一把冰期末期的篝火。(b) $N_{^{14}\text{C}}
= (N_{\text{A}}/12)\times10^{-12} = \qty{5.0e10}{}$; $\lambda =
\qty{3.8e-12}{s^{-1}}$: $\mathcal A \approx \qty{0.2}{Bq}$ ---
每克每分钟约十二次衰变，这正是
各测年实验室的工作计数率。(c) 经过九个半衰期后，放射性活度
下沉至本底以下：钟仍在走，却
再也无法读出。(d) 岩石从未吸入大气中的碳；它们
的计时钟是原生的一类——铀–铅、
钾–氩——半衰期长达数十亿年。
\end{solution}

\begin{solution}{exo:b3:nuclear-physics:8}
(a) $\alpha$ 粒子的 \qty{8.78}{MeV} 远低于
\qty{26}{MeV} 的势垒边缘：经典地说，它既不能逸出，
也不曾进入——然而半衰期竟是 \qty{0.3}{\micro s}。(b)
$\lambda_1 = \ln2/\qty{3e-7}{s} = \qty{2.3e6}{s^{-1}}$;
$\lambda_2 = \ln2/\qty{1.4e17}{s} = \qty{4.9e-18}{s^{-1}}$:
比值 $\approx 5\times10^{23}$。(c) 隧道效应的速率为
$\eu^{-2G}$，其中伽莫夫指数 $\propto Z/\sqrt E$：
将 $E$ 减半会使 $2G$ 增大约五十五个单位，而
$\eu^{55} \approx 10^{24}$——盖革–纳托尔那条离谱的
直线，竟只来自一个指数。(d) 太阳中的质子也
发生隧道效应：核心温度须恰好高到使
麦克斯韦尾部与伽莫夫因子之积足以维持燃烧——
同一道峭壁，从另一侧攀越。
\end{solution}

\begin{solution}{exo:b3:nuclear-physics:9}
(a) 裂变使 ${\sim}235$ 个核子由 \qty{7.6}{} 升至
${\sim}\qty{8.5}{MeV}$ 的结合能：$235\times0.9 \approx
\qty{200}{MeV}$。(b) $N = \qty{2.56e24}{}$ 个原子核：
$E = \qty{8.2e13}{J} \approx 2000$ 吨石油——每
千克。(c) 热功率 \qty{3}{GW} 运行一年，
即 \qty{9.5e16}{J}：约需 \qty{1.2}{t} 的 $^{235}$U。(d) 两个
带正电的碎片在各自的库仑
斥力下弹开，携带约 ${\sim}\qty{170}{MeV}$ 的动能，
在微米量级内转化为热量——反应堆就是一只
以静电反冲为火焰的水壶。
\end{solution}

\begin{solution}{exo:b3:nuclear-physics:10}
(a) 经过 $10^4$ 代：$(1.001)^{10^4} = \eu^{10} \approx
22000$——一秒内功率便从兆瓦级升至数十吉瓦；没有
电动机能如此快地移动控制棒。(b) 当 $k - 1 < \beta$ 时，链
式反应仅凭瞬发中子无法自我繁殖：每一步增长
都要等待约 ${\sim}\qty{10}{s}$ 的缓发中子，因此有效
代时间为零点几秒或更久。(c)
$(1.0005)^{10} \approx 1.005$：每秒约增长百分之零点五——
属于控制棒与操作员的领域。(d) 驱使 $k - 1$ 越过 $\beta$
使反应堆进入 \qty{e-4}{s} 的瞬发时标，
便失去了使反应堆可操控的缓发中子宽限期。
\end{solution}

\begin{solution}{exo:b3:nuclear-physics:11}
(a) $\dot m = L_\odot/c^2 = \qty{4.2e9}{kg/s}$ 为纯质量流率；
氢的消耗率 $\dot m/0.0071 \approx \qty{6e11}{kg/s}$——
即每秒六亿吨。(b) 可燃烧的氢
${\sim}2\times10^{30}\times0.1\times0.75 = \qty{1.5e29}{kg}$:
$t \approx \qty{2.5e17}{s} \approx 8$ 十亿年——太阳
正值中年。(c) 势垒约 ${\sim}\qty{1.4}{MeV}$，
对比热能 $k_{\text{B}}T \approx \qty{1.3}{keV}$，相差千倍
——唯有麦克斯韦分布尾部最快的质子，
借隧道效应（见\cref{exo:b3:nuclear-physics:8}）方能聚变；因此太阳
能稳定燃烧数十亿年，而非爆炸。(d) 其灰烬为
$^4$He——稳定的双幻核，十分安逸；
数个世纪内都没有可供 $\beta$ 衰变的富中子碎片。
\end{solution}

\begin{solution}{exo:b3:nuclear-physics:12}
(a) $\lambda = \ln2/\qty{21600}{s} = \qty{3.2e-5}{s^{-1}}$：
$N = \mathcal A/\lambda = \qty{1.6e13}{}$ 个原子核——
即 \qty{2.6}{ng}。药量以纳克计。(b) 六小时
长于扫描过程，却不及整个周末：剂量自行关断；
一束纯 \qty{140}{keV} 的 $\gamma$ 射线离开人体，交由照相机接收，
不必承受 $\alpha$/$\beta$ 射线灼伤组织的代价。(c) 静止的电子
与正电子湮灭为\emph{两}个
\qty{511}{keV} 的光子，背对背飞出（能量与动量
守恒，见\cref{ch:b3:relativistic-dynamics}）。每一对符合事件
便确定一条穿过肿瘤的直线，成千上万条直线
交汇定位——借守恒律实现的断层成像。(d)
\qty{2}{J/kg} 仅使组织升温半毫开尔文——热学上
不过一口温水；但若以 ${\sim}10^{17}$ 次电离送达，
每千克中，每次都切断 \unit{eV} 量级的化学键，
对分裂中的细胞具有化学致死性。焦耳度量热量；
戈瑞度量的是破坏。
\end{solution}

\begin{solution}{pb:b3:nuclear-physics:1}
\textbf{1.} $\qty{200}{MeV} = \qty{3.2e-11}{J}$：
每秒裂变次数：$20\times10^6/3.2\times10^{-11} = \qty{6.3e17}{}$
次。
\textbf{2.} 每天 $5.4\times10^{22}$ 次裂变 $\times
235/N_{\text{A}}$：约每天消耗 \qty{21}{g} 的 $^{235}$U。
\textbf{3.} 煤：$\qty{1.7e12}{J}/\qty{3e7}{J/kg} \approx 58$
吨每天——质量比为 $\sim3\times10^{6}$，
这\emph{正是}核键与
化学键的 \unit{MeV} 对 \unit{eV} 之比。
\textbf{4.} $235\times(8.5 - 7.6) \approx \qty{210}{MeV}$：
自洽，微小差额由中子与
中微子记账。
\textbf{5.} 碎片动能 $\to$ 电离 $\to$ 热量，
在微米与微秒量级内完成；传导将其交给
水——反应堆正是一只阻停碎片的水壶。
\textbf{6.} 它们继承了铀的 $N/Z \approx 1.55$，远
高于其 $A$ 对应的谷底：每一个都须经历一连串
$\beta^-$ 衰变回落——放射性与生俱来，由
谷地的几何形状所决定。
\textbf{7.} $k$ 为每个中子在下一代中产生的中子数：
$0.999$ 趋于消亡，$1.000$ 保持稳定（即运行中的
反应堆），$1.001$ 增长。
\textbf{8.} 慢中子在原子核附近停留较久，且
$^{235}$U 的裂变截面在热能区急剧增大。
弹性碰撞中，能量向等质量粒子转移最快——
而氢的质量恰与中子相同：水正是
量身定做的慢化剂。
\textbf{9.} 硼-10 吞食中子（$\text{n} + {}^{10}
\text{B} \to \alpha + {}^{7}\text{Li}$）：插入控制棒，中子
被吞食，$k$ 下降；抽出控制棒，$k$ 上升——这便是节流阀。
\textbf{10.} $(1.0005)^{10^4} = \eu^{5} \approx 150$：功率每秒增长
约一百五十倍——机械根本无从应对。
\textbf{11.} 在 \qty{0.1}{s} 的缓发中子时标上：
$(1.0005)^{10} \approx 1.005$，每秒约增长百分之零点五。规则：
须使 $|k - 1|$ 远小于 $\beta = 0.007$，从而使缓发
中子始终握有关键一票。
\textbf{12.} 沸腾移除了慢化剂：中子保持快速，
裂变减少，$k$ 下降——这种水慢化设计
自行节流（一种负反馈，其石墨慢化的同类
反应堆所缺乏）。
\textbf{13.} 衰变热：碎片的放射性——
第一部分第 6 题——遵循半衰期，而不受控制棒影响，且
停堆后不久仍能释放兆瓦级功率。
\textbf{14.} 在水中超过光速：$v > c/n$。
\textbf{15.} $v > 0.752c$：$\gamma = 1.51$，故 $E_{\text{k}} >
0.51\times\qty{511}{keV} \approx \qty{0.26}{MeV}$.
\textbf{16.} 碎片衰变产生的 $\beta$ 电子携带
\unit{MeV} 量级的能量，堆芯的 $\gamma$ 射线借康普顿效应
将电子打到相近能量：二者均轻松越过 \qty{0.26}{MeV}——水池
中充满了合格的电子。
\textbf{17.} 切伦科夫光谱向短
波长方向增强（强度 $\propto1/\lambda^2$）：人眼接收的是
蓝光与紫光——颜色正由光谱的斜率决定。
\textbf{18.} 水对 $\gamma$ 射线与中子的衰减
按指数规律；八米约合数十个半值层，因此
阳台处于本底剂量。辉光正是水
在\emph{吸收}辐射能量——可见的证据表明
屏蔽层正在吞噬它。
\textbf{19.} 短寿命碎片数分钟内即衰变殆尽
（辉光渐暗）；长寿命尾部——即第 13 题所述的
同一衰变热——仍保持微弱辉光与实际热负荷，持续
数日。
\textbf{20.} 取 $T_{1/2} = \qty{66}{h}$，母核每周
约减少 ${\sim}6$ 倍：库存自行衰减殆尽，
因此医院全靠每周一次的“锝牛”度日，它由
与本堆类似的反应堆送来。
\textbf{21.} $96/66 = 1.45$ 个半衰期：$2^{-1.45} \approx
0.36$——约剩 \qty{36}{GBq}。
\textbf{22.} 热化至室温后，$\lambda =
h/\sqrt{3m_{\text{n}}k_{\text{B}}T} \approx \qty{1.5}{\angstrom}$
（见\cref{exo:b3:crystalline-solids:5}）：生来便与晶格
间距相匹配。
\textbf{23.} 五年即可使半衰期在数月以内的核素全部衰减殆尽：
放射性活度与衰变热下降数个数量级，直至该
元素可装入屏蔽容器运输，而不必再泡在水池中。
\textbf{24.} ${\sim}\qty{200}{kW}$——相当于一百台家用
电暖器加热一座百吨级水池：每天至多
升温数十开尔文，最坏情形下自然对流即可轻松带走——非能动
安全，全凭热容之大。
\textbf{25.} $B/A$ 每次裂变释放 \qty{200}{MeV}，每天
消耗 \qty{21}{g}；缓发中子提供使 $k$
可操控的秒级时间；水兼具慢化、冷却与屏蔽之用——并
因其正在运行而发出蓝光；而最关键的产物
则装入小瓶，供周一的医院使用。
\end{solution}
